\chapter{Requirement Analysis}
\setcounter{equation}{0}

\section{Software Requirements}
BrowseBlock requirements are derived from the Software Requirements Specification.
The system consists of a Chromium Manifest V3 extension and a React web dashboard
supported by a FastAPI service. Mandatory behavior is expressed as a system
obligation, while optional behavior depends on browser permissions, available
evidence, and the selected deployment configuration.

\section{User Classes and Characteristics}
General users require short warnings, a clear recommendation, and a way to continue
or return to browsing. Authorized users of the dashboard require scan forms, result
history, threat evidence, and analytics. Security analysts require detailed
indicators and report context for investigation. Administrators require controlled
configuration and access management. These classes differ in privilege, but none
should receive model credentials, database credentials, or raw data beyond the
authorized workflow.

\subsection{Functional Requirements}
The system shall authenticate dashboard users through Firebase Authentication and
shall verify identity at the FastAPI boundary. It shall accept supported URLs, QR
codes, emails, webpages, and downloaded files for analysis. It shall select the
appropriate detector from the triggering browser event or the user-selected scan
type, validate input, and return a common result containing a risk level, evidence,
analysis completeness, and recommended or completed protective action.

The extension shall monitor supported navigation, webpage scripts and DOM behavior,
redirects, frames, notification prompts, downloads, and observable extension
activity. It shall provide a warning or restriction when the browser APIs,
permissions, evidence, and policy support that action. The dashboard shall allow
authorized users to submit manual scans, review reports and analytics, and inspect
the evidence behind a decision. The system shall never require a real password as
an input to a detection workflow.

\subsubsection{Functional Requirement Summary}
The main functional requirements of BrowseBlock are:
\begin{itemize}
\item Authenticate dashboard users and authorize access to scans, reports, and analytics.
\item Monitor browser navigation, webpages, redirects, frames, notifications, downloads, and extension activity.
\item Detect URL and QR-code phishing using URL normalization, feature extraction, and classification.
\item Detect reflected, stored, and DOM-based cross-site scripting indicators.
\item Detect credential-harvesting forms, suspicious domains, scripts, and scam pop-ups without collecting passwords.
\item Track redirect chains, evaluate final destinations, and warn about malicious navigation.
\item Detect clickjacking through suspicious frames, overlays, positioning, and interaction patterns.
\item Detect cryptojacking through mining scripts, WebAssembly, workers, mining endpoints, and resource behavior.
\item Analyze downloaded executable, PDF, DOC, DOCX, and ZIP files using static inspection, hashes, rules, and models.
\item Classify suspicious files and support warning, cancellation, quarantine, or release according to policy.
\item Detect email phishing using sender details, links, authentication signals, machine learning, retrieval, and semantic analysis.
\item Detect notification scams using origin, prompt text, urgency, destinations, and redirects.
\item Detect browser-extension abuse through permissions, host access, injected scripts, API use, and page changes.
\item Display Safe, Suspicious, or High-risk results with evidence, analysis status, and recommended protective action.
\item Record detection events and clearly report unsupported inputs, unavailable evidence, failed analysis, and incomplete results.
\end{itemize}

\subsection{Non-Functional Requirements}
\begin{itemize}
\item \textbf{Security:} protected requests shall use HTTPS in deployment, Firebase
ID tokens shall be verified server-side, and secrets shall not be exposed to the
extension or dashboard.
\item \textbf{Privacy:} collection shall be limited to data required for detection
and authorized reporting. Email, webpage, and file content shall be submitted for
deeper analysis only when permitted by the workflow.
\item \textbf{Performance:} lightweight checks shall remain event-driven, while
expensive models and retrieval workflows shall execute in the backend.
\item \textbf{Reliability:} timeouts, unsupported formats, failed parsers, and
unavailable browser events shall produce an explicit incomplete-analysis state,
not a false Safe verdict.
\item \textbf{Usability:} warnings shall present the risk, reason, evidence, and
available action in language suitable for non-specialist users.
\item \textbf{Maintainability:} detectors, models, rules, and user interfaces shall
remain modular and independently testable.
\end{itemize}

\section{Design and Implementation Constraints}
The extension must operate within Manifest V3 lifecycle and permission rules.
Content scripts can inspect only the pages and events exposed by the browser, and
some protected or cross-origin content may be unavailable. Backend analysis depends
on network connectivity, model availability, file-size limits, parser safety, and
the response time of external reputation or language services. The system shall
therefore use timeouts, input-size limits, isolated processing, and explicit
fallback states rather than silently treating a failed operation as Safe.

\section{External Interface Requirements}
\subsection{User Interfaces}
The extension shall provide concise warnings, scan status, risk classification,
indicator summaries, and an available action such as continue, cancel, or review.
The dashboard shall provide authentication, URL and QR scanning, email and file
submission, report viewing, analytics, and help describing analysis limitations.

\subsection{Software Interfaces}
The extension and dashboard communicate with REST-style FastAPI endpoints. The
backend validates structured requests, invokes the relevant detector, and returns a
consistent JSON result. Firebase Authentication supplies user identity; client
components do not connect directly to protected backend resources or model files.

\subsection{Operating Environment}
The extension targets supported Chromium browsers that provide the permissions
needed for the selected monitoring features. The dashboard runs in a modern browser.
The backend runs in a Python environment with FastAPI and the configured model,
parsing, QR-decoding, and language-analysis dependencies. A stable network is
required for centralized analysis and authentication.

\subsection{Communication and Data Interfaces}
Extension and dashboard requests shall use structured JSON over HTTPS in deployment.
Protected requests shall carry a Firebase ID token, while the backend shall validate
the token and determine the authorized user or role. A scan request identifies the
input type, source, and permitted content. A result includes the risk level,
confidence or supporting evidence, analysis status, timestamps, and actual action.
MongoDB stores authorized application data through backend operations; clients do
not connect directly to the database.

\section{System Features}
\subsection{URL and Fake QR Phishing}
The URL detector parses protocol, host, domain, subdomains, path, query, length,
symbols, IP usage, and suspicious keywords. XGBoost classifies the extracted
features. A QR workflow decodes the supported image, verifies whether its payload
is a URL, and passes that URL to the same detector.

The stimulus is a supported navigation event or a QR code submitted through the
dashboard. The response sequence is capture or decode, normalize, feature
extraction, XGBoost classification, explanation, and warning or restriction where
supported. Malformed, unusually long, or encoded input must be handled without
crashing. The detector shall be evaluated using precision, recall, and
false-positive measurements while minimizing storage of sensitive URL parameters.

\subsection{Cross-Site Scripting}
The system inspects supported scripts, source-to-sink patterns, DOM changes, and
observable requests for reflected, stored, and DOM-based XSS indicators. It reports
indicators and actual enforcement separately because detection of a suspicious
pattern does not prove that script execution was prevented.

The detector shall examine suspicious sources and sinks, inline scripts, event
handlers, dynamically created script elements, unexpected script sources, and
observable DOM manipulation. It shall distinguish reflected, stored, and DOM-based
indicators where evidence permits. Conclusions remain limited to content scripts
and browser APIs visible at runtime, and legitimate page behavior must not be
disrupted unnecessarily.

\subsection{Credential Harvesting and Redirects}
Credential-harvesting analysis identifies deceptive forms, password fields,
mismatched page identity, external submission destinations, and suspicious scripts
without collecting real credentials. Redirect analysis evaluates navigation chains,
delayed JavaScript or HTML redirects, intermediary domains, frames, and the final
destination where those events are visible to the browser.

For credential harvesting, the stimulus is a login or sensitive-data form. The
system compares the displayed identity, registered domain, form action, password
fields, external scripts, and submission destination without collecting a real
password. For redirects, it records observable source and destination URLs,
evaluates the final destination using URL-risk information, avoids repeated
analysis of the same event, and reports when the browser did not expose a complete
chain.

\subsection{Clickjacking and Cryptojacking}
Clickjacking analysis correlates invisible or deceptive overlays, iframe visibility,
stacking, positioning, and interaction patterns. Cryptojacking analysis examines
scripts, WebAssembly, workers, mining endpoints, and abnormal resource indicators.
Legitimate overlays and resource-intensive pages require multiple indicators before
a high-risk decision.

Clickjacking processing begins when a page contains frames or embedded interactive
content. It examines iframe attributes, visibility, position, overlay relationships,
stacking, and interaction structure, then warns or restricts only when a reliable
browser-supported mechanism exists. Cryptojacking processing begins when scripts,
WebAssembly, workers, or mining-related network behavior are observed. It combines
configured mining-resource indicators with available resource evidence while
minimizing monitoring overhead and false positives on legitimate intensive pages.

\subsection{Malicious File Downloads}
Supported downloads, including ZIP, PDF, EXE, DOC, and DOCX files, are identified
from their content and extension. The workflow can calculate SHA-256 hashes, apply
format-specific static analysis and deterministic rules, use reputation information,
and invoke the configured model. Untrusted files are not executed merely for
analysis.

When a supported download starts, the extension identifies the file and transfers
only the required content or metadata to the controlled workflow. The backend
applies file-specific parsers, static rules, SHA-256 and reputation checks, and
XGBoost where configured. A native helper may provide controlled file access,
quarantine, and release support. Size, archive, time, and resource limits shall
be enforced, and the final record shall distinguish allow, warn, cancel, and
quarantine actions.

\subsection{Email Phishing}
The extension performs lightweight screening using Random Forest features and
available SPF, DKIM, and DMARC signals. Authorized dashboard analysis can extract
URLs, retrieve relevant knowledge through RAG, and use an LLM to assess urgency,
impersonation, social engineering, and credential harvesting. Extracted URLs are
also evaluated by the URL detector.

The first stage screens sender and domain differences, link-text mismatch,
IP-based URLs, suspicious terms, link count, script indicators, and available
authentication signals. A dashboard investigation is a second stage: after user
authorization, RAG retrieves context, links and patterns are extracted, the URL
detector assesses extracted links, and the LLM summarizes semantic indicators.
SPF, DKIM, and DMARC success is evidence only, not proof of safety. Retrieval,
inference, and external requests require configurable timeouts, and failed analysis
must remain visibly incomplete.

\subsection{Browser Notification Scams and Extension Abuse}
Notification analysis evaluates origin, text, urgency, fake security warnings,
reward or payment language, destinations, and redirects. Extension-abuse analysis
compares declared permissions, host access, observable API use, injected scripts,
and page changes. Excessive permission alone is treated as an indicator rather
than proof of malicious behavior.

Notification analysis begins when a page requests permission or displays a related
prompt. It correlates origin, prompt text, associated URL, fake security or prize
language, and post-click behavior while minimizing interference with legitimate
notifications. Extension-abuse analysis examines observable permissions,
unexpected webpage tampering, suspicious injection, and host access. If browser
APIs do not expose enough information, the system shall state that limitation and
avoid claiming a complete assessment.

\section{Risk Classification}
Every detector maps its evidence to one of three user-facing levels: Safe,
Suspicious, or High-risk. The result also records the detector category, evidence
indicators, timestamp, analysis completeness, and protective action. A completed
scan is not a guarantee that an item is safe, and the report must identify
unsupported content, incomplete visibility, and unavailable enforcement.

\section{Logging, Audit, and Reporting Requirements}
The system shall record sufficient metadata to explain a result and troubleshoot
the workflow: scan type, detector, time, status, indicators, risk level, and
protective-action outcome. Sensitive URL parameters, email bodies, webpage content,
and file content shall not be retained unnecessarily. Audit records shall distinguish
a requested action from a completed action and shall be protected by authentication,
authorization, and the configured retention policy.

\section{User Documentation and Future Requirements}
Documentation shall explain installation, supported browser versions, supported
file types and limits, scan workflows, report interpretation, risk decisions,
privacy behavior, and troubleshooting for failed or unsupported operations. The
modular architecture shall support future threat categories, updated models, new
browser APIs, improved threat intelligence, and additional dashboard analytics.
The implementation must finalize the browser/API matrix, model versions and
datasets, external providers, file limits, retention periods, deployment
environment, and warning, cancellation, quarantine, and release policies.

\section{Proposed Solution}
The proposed solution is an event-driven extension backed by a centralized
FastAPI service and an authenticated React dashboard. The extension performs
low-latency collection and user feedback. FastAPI validates input, verifies
identity, dispatches detectors, coordinates XGBoost, Random Forest, static rules,
QR decoding, file analysis, reputation services, and RAG/LLM workflows, and writes
authorized records to MongoDB. The dashboard provides deliberate scans and deeper
explanations without bypassing the backend security boundary.

\section{Scope and Limitations}
BrowseBlock is an additional browser-security layer, not a replacement for
antivirus, browser security controls, or enterprise monitoring. Detection is
bounded by Chromium permissions, browser API visibility, network availability,
model quality, external-service availability, and the evidence observable at
runtime. The first prototype will prioritize representative modules and controlled
evaluation rather than claiming complete prevention of every attack.
